\subsubsection*{Orientation}

\begin{remark*}[Section map]
This section uses the preceding syntax, semantics, and algebra of propositions
as its starting point. It first isolates literals and clauses, then introduces
negation normal form, disjunctive normal form, conjunctive normal form,
canonical normal forms, and semantic classification uses. Later arguments
will use these forms to organize satisfiability questions, proof search,
functional completeness, and resolution-style methods.
\end{remark*}

\begin{remark*}[Why normal forms matter]
Normal forms trade arbitrary expression for regular structure. In a truth-table
argument, regular structure helps classify formulas as tautologies,
contradictions, satisfiable formulas, or contingencies. In proof search and
satisfiability procedures, regular structure turns a formula into a finite
combinatorial object that can be searched. In functional completeness, normal
forms show how complicated truth functions can be built from simpler
connectives.
\end{remark*}
